\subsection{LIE ALGEBRA OF THE GROUP OF INVERTIBLE ELEMENTS OF AN ALGEBRA}

Let A be a complete normable associative algebra with unit element \(e\) . Let \(\mathbf{A}^*\) be the group of invertible elements of A. We have seen (§ 1, no. 1) that \(\mathbf{A}^*\) is an open submanifold of A and is a Lie group. Let \(\mathbf{G}\) be a Lie group and \(f\) a morphism of the Lie group \(\mathbf{G}\) into the Lie group \(\mathbf{A}^*\) , \(f\) can be considered as an analytic mapping of \(\mathbf{G}\) into the complete normable space A. Hence, if \(t \in \mathcal{T}^{(\infty)}(\mathbf{G})\) , we can form \(\langle t, f \rangle\) , which is an element of A.

PROPOSITION 33. The mapping \(t \mapsto \langle t, f \rangle\) is a morphism of the algebra \(\mathcal{T}^{(\infty)}(G)\) into the algebra A.

It suffices to verify that, if \(t\) and \(t'\) are point distributions on \(G\) , then \(\langle t * t', f \rangle = \langle t, f \rangle \langle t', f \rangle\) . But

\[
\begin{array}{r l} \langle t * t ^ {\prime}, f \rangle & = \langle t \otimes t ^ {\prime}, (g, g ^ {\prime}) \mapsto f (g g ^ {\prime}) \rangle \\ & = \langle t \otimes t ^ {\prime}, (g, g ^ {\prime}) \mapsto f (g) f (g ^ {\prime}) \rangle \\ & = \langle t, f \rangle \langle t ^ {\prime}, f \rangle \\ & \quad (Diff. \& Anal. Man., R, 13.4.3). \end{array}
\]

The morphism of Proposition 33 is said to be associated with f.

Take G to be the group A* itself and f to be the identity mapping \(\iota\) of A*. We obtain a morphism, called canonical, of the algebra \(\mathcal{T}^{(\infty)}(\mathbf{A}^*)\) into the algebra A. The tangent space \(\mathrm{T}_{e}(\mathbf{A}^*)\) is canonically identified with A; and if \(t \in \mathrm{T}_{e}(\mathbf{A}^*)\) , the definition of this identification is such that \(\langle t, \iota \rangle = t\) . Then Proposition 33 implies the following corollary:

COROLLARY. The canonical mapping \(\zeta\) of \(\mathbf{L}(\mathbf{A}^{*})\) into \(\mathbf{A}\) is an isomorphism of the Lie algebra, \(\mathbf{L}(\mathbf{A}^{*})\) onto the Lie algebra \(\mathbf{A}\) . In other words,

\[
\zeta ([ a, b ]) = \zeta (a) \zeta (b) - \zeta (b) \zeta (a)
\]

for all \(a, b\) in \(\mathrm{L}(\mathbf{A}^*)\) . If \(\mathbf{K}\) is of characteristic \(p > 0\) , then \(\zeta(a^p) = \zeta(a)^p\) for all \(a \in \mathrm{L}(\mathbf{A}^*)\) .

Henceforth L(A*) and A are identified by means of the isomorphism ζ.

The canonical morphism of \(\mathcal{T}^{(\infty)}(\mathbf{A}^*)\) into \(\mathbf{A}\) has been obtained as a special case of the morphism of Proposition 33. But it is possible to argue in the opposite direction:

PROPOSITION 34. Let H be a Lie group, A a unital complete normable associative algebra and \(\phi: \mathrm{H} \to \mathrm{A}^*\) a Lie group morphism. The associated morphism \(\phi'\) of \(\mathcal{T}^{(\infty)}(\mathrm{H})\) into A is obtained by composing \(\phi_*\) with the canonical morphism of \(\mathcal{T}^{(\infty)}(\mathrm{A}^*)\) into A. In particular, \(\phi'(x) = \mathrm{L}(\phi)(x)\) for all \(x \in \mathrm{L}(\mathrm{H})\) .

Let \(i\) be the identity mapping of \(\mathbf{A}^*\) into A. Then, for all \(t\in \mathcal{T}^{(\infty)}(\mathrm{H})\) ,

\[
\begin{array}{r l} \phi^ {\prime} (t) = \langle t, \phi \rangle & = \langle t, i \circ \phi \rangle \\ & = \langle \phi_ {*} (t), i \rangle \quad (Diff. \& Anal. Man., R, 13.2.3). \end{array}
\]

